%============================================================
%  Bd. 16 - Wohlordnungen und Auswahlaxiom
%============================================================

\documentclass[main.tex]{subfiles}

\ifSubfilesClassLoaded{
  \BandSetupStandalone{B16}
}{
}

\title{Bd. 16 - Wohlordnungen und Auswahlaxiom}
\author{Martin Kunze}
\date{}
\setcounter{file}{16}

\hypersetup{
  pdftitle={Bd. 16 - Wohlordnungen und Auswahlaxiom},
  pdfauthor={Martin Kunze}
}

\begin{document}

\title{Bd. 16 - Wohlordnungen und Auswahlaxiom}
\author{Martin Kunze}
\date{}
\hypersetup{pageanchor=false}
\maketitle
\hypersetup{pageanchor=true}
\setcounter{tocdepth}{3}
\tableofcontents
\setcounter{file}{16}
\renewcommand{\FormulaBandID}{16}
\chapter{Einleitung}

Wohlordnungen sind totale Ordnungen, in denen jede nichtleere Teilmenge ein
kleinstes Element besitzt. Dieser Band entwickelt Wohlordnung und
Wohlordenbarkeit, formuliert den Wohlordnungssatz und beweist seine Äquivalenz
mit dem Auswahlaxiom. Die natürlichen Zahlen dienen als grundlegendes Beispiel.

\chapter{Wohlordnungen}
\section{Wohlordnungen, Wohlordnungssatz und Auswahlaxiom}
\subsection{Axiom einer Wohlordnung}

\FormulaAxiomDeltaK[Wohlordnungsprinzip]{%
  T\subseteq A \dsep T\neq\varnothing
  \vdash \exists m\in T\,\forall x\in T\,\bigl(m\leq x\bigr)%
}{WellOrderPrinciple}{
  \DeltaRow{Mengen}{A\dsep T \dsep m \dsep x}
  \DeltaRow{Relationsoperatoren}{\leq}
}

\subsection{Definition der Wohlordnung}

\FormulaDefDeltaK[Begriff der Wohlordnung]{\WellOrd{A,\leq}}{Wohlordnung}{
  \DeltaRow{Mengen}{A\dsep T \dsep m \dsep x \dsep y \dsep z}
  \DeltaRow{Relationsoperatoren}{\leq}
  \DeltaRow{\textbf{Axiome}}{}
  \DeltaRow{Totale Ordnung}
           {\TotOrd{A,\leq}}
           [\FormulaRefAuto{Totale Ordnung}]
  \DeltaRow{Wohlordnungsprinzip}{%
    \begin{aligned}
      &T\subseteq A \dsep T\neq\varnothing\\
      &\vdash \exists m\in T\,\forall x\in T\,\bigl(m\leq x\bigr)
    \end{aligned}}
           [\FormulaRefAuto{WellOrderPrinciple}[ax]]
  \DeltaRow{\textbf{Neue Symbole}}{}
  \DeltaPrem{Wohlordnung}{\WellOrd{A,\leq}}
}

\FormulaDefDeltaK[Begriff der Wohlordenbarkeit]{%
  \WellOrdable{A}\coloneqq
  \exists\,\preccurlyeq\;\WellOrd{A,\preccurlyeq}%
}{WellOrdable}{
  \DeltaRow{Mengen}{A}
  \DeltaRow{Relationsoperatoren}{\preccurlyeq}
  \DeltaRow{Wohlordnung}
           {\WellOrd{A,\preccurlyeq}}
           [\FormulaRefAuto{Wohlordnung}]
  \DeltaRow{\textbf{Neue Symbole}}{}
  \DeltaPrem{Wohlordenbare Menge}{\WellOrdable{A}}
}

\subsection{Wohlordnungssatz}

\FormulaDefDeltaK[Wohlordnungssatz]{%
  \WOP\coloneqq \forall A\,\WellOrdable{A}%
}{WOP}{
  \DeltaRow{Mengen}{A}
  \DeltaRow{Wohlordenbarkeit}
           {\WellOrdable{A}}
           [\FormulaRefAuto{WellOrdable}]
}
\subsection{Auswahlaxiom und Äquivalenz zum Wohlordnungssatz}

\FormulaDefDeltaK[Auswahlrelation auf nichtleeren Teilmengen]{%
  \ChoiceRel{A}\coloneqq
  \{(X,a)\in \PowNE{A}\times A\mid a\in X\}%
}{ChoiceRelation}{
  \DeltaDecl{Mengen}{A \dsep X \dsep a}
  \DeltaDecl{Relation}{\ChoiceRel{A}}
}

\FormulaDefDeltaK[Auswahlfunktion auf nichtleeren Teilmengen]{%
  \ChoicePow{C}{A}\coloneqq
  C\colon \PowNE{A}\to A
  \land \forall X\in \PowNE{A}\,\bigl(C(X)\in X\bigr)%
}{ChoicePow}{
  \DeltaDecl{Mengen}{A \dsep X}
  \DeltaDecl{Funktion}{C}
  \DeltaRow{Definitionsbereich}{\PowNE{A}}
  \DeltaRow{Auswahlbedingung}
           {\forall X\in \PowNE{A}\,\bigl(C(X)\in X\bigr)}
}

\FormulaThmDeltaK[Auswahlrelation ist total]{%
  \TotRel{\ChoiceRel{A},\PowNE{A},A}%
}{ChoiceRelTotal}{
  \DeltaDecl{Mengen}{A \dsep X \dsep a}
  \DeltaDecl{Relation}{\ChoiceRel{A}}
}
\begin{tabproof}
  \proofstep{}{%
    \ChoiceRel{A}\subseteq \PowNE{A}\times A%
  }{%
    \FormulaRefAuto{A=\{ x \in B \mid P(x)\}\vdash  A\subseteq B}{%
      \FormulaRefAuto{\ChoiceRel{A}\coloneqq
      \{(X,a)\in \PowNE{A}\times A\mid a\in X\}}}%
  }

  \proofstep{2}{X\in \PowNE{A}}{\rA}
  \proofstep{2}{X\in \powerset(A)}{%
    \FormulaRefAuto{x \in A \setminus B \vdash x \in A}{2}}
  \proofstep{2}{X\notin \{\varnothing\}}{%
    \FormulaRefAuto{x\in A\setminus B\vdash x\notin B}{2}}
  \proofstep{2}{X\neq\varnothing}{%
    \FormulaRefAuto{x \notin \{a\} \eqvdash x \neq a}{4}}
  \proofstep{2}{X\subseteq A}{%
    \FormulaRefAuto{x \in \powerset(A)\eqvdash x \subseteq A}{3}}
  \proofstep{2}{\exists a\,(a\in X)}{%
    \FormulaRefAuto{A\neq\varnothing \eqvdash \exists x\,(x \in A)}{5}}

  \proofstep{8}{a\in X}{\rA}
  \proofstep{2,8}{a\in A}{%
    \FormulaRefAuto{A\subseteq B,\, x\in A \vdash x\in B}{6,8}}
  \proofstep{2,8}{(X,a)\in \PowNE{A}\times A}{%
    \FormulaRefAuto{a\in A\dsep b\in B\vdash (a,b)\in A\times B}{2,9}}
  \proofstep{2,8}{(X,a)\in \ChoiceRel{A}}{%
    \rIE{\FormulaRefAuto{\ChoiceRel{A}\coloneqq
      \{(X,a)\in \PowNE{A}\times A\mid a\in X\}},
      \FormulaRefAuto{x \in A\dsep P(x)\vdash x \in \{x \in A \mid P(x)\}}{10,8}}}
  \proofstep{2,8}{\exists a\,\bigl((X,a)\in \ChoiceRel{A}\bigr)}{\rEI{11}}
  \proofstep{2}{\exists a\,\bigl((X,a)\in \ChoiceRel{A}\bigr)}{\rEE{7,8,12}}

  \proofstep{}{%
    X\in \PowNE{A}\rightarrow \exists a\,\bigl((X,a)\in \ChoiceRel{A}\bigr)%
  }{\rRI{2,13}}
  \proofstep{}{%
    \TotRel{\ChoiceRel{A},\PowNE{A},A}%
  }{\FormulaRefAuto{Totale Relation}{1,14}}
\end{tabproof}

\FormulaThmDeltaK[Elemente der Auswahlrelation]{%
  (X,a)\in \ChoiceRel{A}\vdash a\in X%
}{ChoiceRelElim}{
  \DeltaDecl{Mengen}{A \dsep X \dsep a}
  \DeltaDecl{Relation}{\ChoiceRel{A}}
}
\begin{tabproof}
  \proofstep{1}{(X,a)\in \ChoiceRel{A}}{\rA}
  \proofstep{1}{%
    (X,a)\in \{(X,a)\in \PowNE{A}\times A\mid a\in X\}%
  }{%
    \rIE{\FormulaRefAuto{\ChoiceRel{A}\coloneqq
      \{(X,a)\in \PowNE{A}\times A\mid a\in X\}},1}}
  \proofstep{1}{(X,a)\in \PowNE{A}\times A\land a\in X}{%
    \FormulaRefAuto{x \in \{u \in A \mid P(u)\} \eqvdash x \in A \land P(x)}{2}}
  \proofstep{1}{a\in X}{\rAEb{3}}
\end{tabproof}

\FormulaThmDeltaK[Auswahlrelation liefert Auswahlfunktion]{%
  \begin{aligned}[t]
    &\exists C\colon \PowNE{A}\to A\,
      \forall X\in \PowNE{A}\,
        \bigl((X,C(X))\in \ChoiceRel{A}\bigr)\vdash{}\\
    &\exists C\,\ChoicePow{C}{A}
  \end{aligned}%
}{ChoiceRelToChoicePow}{
  \DeltaDecl{Mengen}{A \dsep X}
  \DeltaDecl{Funktionen}{C}
  \DeltaDecl{Relation}{\ChoiceRel{A}}
}
\begin{tabproofwide}
  \proofstepwidestar[1]{%
    \exists C\colon\PowNE{A}\to A\,
      \forall X\in\PowNE{A}\,
        \bigl((X,C(X))\in\ChoiceRel{A}\bigr)}{\rA}
  \proofstepwidestar[2]{%
    C\colon\PowNE{A}\to A\land
      \forall X\in\PowNE{A}\,
        \bigl((X,C(X))\in\ChoiceRel{A}\bigr)}{\rA}
  \proofstepwidestar[2]{C\colon\PowNE{A}\to A}{\rAEa{2}}
  \proofstepwidestar[2]{%
    \forall X\in\PowNE{A}\,
      \bigl((X,C(X))\in\ChoiceRel{A}\bigr)}{\rAEb{2}}
  \proofstepwidestar[5]{X\in\PowNE{A}}{\rA}
  \proofstepwidestar[2,5]{(X,C(X))\in\ChoiceRel{A}}{%
    \rRE{5,\rUE{4}}}
  \proofstepwidestar[2,5]{C(X)\in X}{%
    \FormulaRefAuto{ChoiceRelElim}[thm]{6}}
  \proofstepwidestar[2]{%
    \forall X\in\PowNE{A}\,\bigl(C(X)\in X\bigr)}{%
    \rUI{\rRI{5,7}}}
  \proofstepwidestar[2]{\ChoicePow{C}{A}}{%
    \FormulaRefAuto{ChoicePow}[def]{\rAI{3,8}}}
  \proofstepwidestar[2]{\exists C\,\ChoicePow{C}{A}}{\rEI{9}}
  \proofstepwidestar[1]{\exists C\,\ChoicePow{C}{A}}{\rEE{1,2,10}}
\end{tabproofwide}

\FormulaThmDeltaK[Surjektionsform liefert Auswahlfunktion auf nichtleeren Teilmengen]{%
  \ACsur\vdash \exists C\,\ChoicePow{C}{A}%
}{ACsurChoicePow}{
  \DeltaDecl{Mengen}{A \dsep X}
  \DeltaDecl{Funktionen}{C \dsep s}
  \DeltaDecl{Relation}{\ChoiceRel{A}}
}
\begin{tabproof}
  \proofstep{1}{\ACsur}{\rA}
  \proofstep{1}{%
    \forall A\,\forall B\,\forall G\,
    \bigl(G\colon B\sur A\rightarrow
    \exists F\colon A\to B\,G\circ F=\Id_A\bigr)%
  }{\rIE{\FormulaRefAuto{ACsur},1}}
  \proofstep{}{\TotRel{\ChoiceRel{A},\PowNE{A},A}}{\FormulaRefAuto{ChoiceRelTotal}}
  \proofstep{}{\pi_1\restriction_{\ChoiceRel{A}}\colon \ChoiceRel{A}\sur \PowNE{A}}{%
    \FormulaRefAuto{\pi_1\restriction_R\colon R\sur A}{3}}
  \proofstep{1}{%
    \exists s\colon \PowNE{A}\to \ChoiceRel{A}\,
    \bigl(\pi_1\restriction_{\ChoiceRel{A}}\circ s=\Id_{\PowNE{A}}\bigr)%
  }{\rRE{\rUE{2},4}}
  \proofstep{1}{%
    \exists C\colon \PowNE{A}\to A\,
    \forall X\in \PowNE{A}\,\bigl((X,C(X))\in \ChoiceRel{A}\bigr)%
  }{%
    \FormulaRefAuto{\exists s\colon A\to R\,\bigl(\pi_1\restriction_R\circ s = \Id_A\bigr)
    \vdash \exists F\colon A\to B\,\forall x\in A\,\bigl((x,F(x))\in R\bigr)}{5}}
  \proofstep{1}{\exists C\,\ChoicePow{C}{A}}{\FormulaRefAuto{ChoiceRelToChoicePow}{6}}
\end{tabproof}

Die für Zermelos Konstruktion benötigte transfinite Rekursion wurde bislang
nicht entwickelt. Das folgende Konstruktionsprinzip kapselt deshalb genau
diesen noch nicht formalisierten Schritt.

\FormulaAxiomDeltaK[Zermelos Konstruktionsprinzip]{%
  \ChoicePow{C}{A}\vdash
  \exists\,\preccurlyeq\;\WellOrd{A,\preccurlyeq}%
}{ZermeloConstruction}{
  \DeltaDecl{Mengen}{A}
  \DeltaDecl{Funktion}{C}
  \DeltaRow{Relationsoperatoren}{\preccurlyeq}
  \DeltaRow{Auswahlfunktion}
           {\ChoicePow{C}{A}}
           [\FormulaRefAuto{ChoicePow}]
  \DeltaRow{Konstruktionsresultat}
           {\exists\,\preccurlyeq\;\WellOrd{A,\preccurlyeq}}
}

\FormulaThmDeltaK[Zermelos Wohlordnungskonstruktion]{%
  \exists C\,\ChoicePow{C}{A}\vdash \WellOrdable{A}%
}{ZermeloWO}{
  \DeltaDecl{Mengen}{A}
  \DeltaDecl{Funktion}{C}
}
\begin{tabproof}
  \proofstep{1}{\exists C\,\ChoicePow{C}{A}}{\rA}
  \proofstep{2}{\ChoicePow{C}{A}}{\rA}
  \proofstep{2}{\exists\,\preccurlyeq\;\WellOrd{A,\preccurlyeq}}{%
    \FormulaRefAuto{ZermeloConstruction}{2}}
  \proofstep{2}{\WellOrdable{A}}{\rIE{\FormulaRefAuto{WellOrdable},3}}
  \proofstep{1}{\WellOrdable{A}}{\rEE{1,2,4}}
\end{tabproof}

\FormulaThmDeltaK[Auswahlaxiom impliziert Wohlordnungssatz]{%
  \ACsur\vdash \WOP%
}{ACsurToWOP}{
  \DeltaDecl{Mengen}{A}
}
\begin{tabproof}
  \proofstep{1}{\ACsur}{\rA}
  \proofstep{1}{\exists C\,\ChoicePow{C}{A}}{\FormulaRefAuto{ACsurChoicePow}{1}}
  \proofstep{1}{\WellOrdable{A}}{\FormulaRefAuto{ZermeloWO}{2}}
  \proofstep{1}{\forall A\,\WellOrdable{A}}{\rUI{3}}
  \proofstep{1}{\WOP}{\rIE{\FormulaRefAuto{WOP},4}}
\end{tabproof}

\FormulaDefDeltaK[Faserminimenge zu einer Wohlordnung]{%
  L_{G,\preccurlyeq}\coloneqq
  \{\,b\in B\mid
    \forall c\in B\,
      \bigl(G(c)=G(b)\rightarrow b\preccurlyeq c\bigr)\,\}%
}{FiberMinSet}{
  \DeltaDecl{Mengen}{A \dsep B\dsep b \dsep c}
  \DeltaDecl{Funktion}{G}
  \DeltaRow{Relationsoperatoren}{\preccurlyeq}
  \DeltaDecl{Mengen}{L_{G,\preccurlyeq}}
}

\FormulaThmDeltaK[Wohlgeordnete Fasern liefern eine Auswahlmenge]{%
  \WellOrd{B,\preccurlyeq}\dsep G\colon B\sur A
  \vdash
  \forall X\in \Fib_G[A]\,
    \exists! b\,(b\in X\land b\in L_{G,\preccurlyeq})%
}{WellOrdFiberChoiceSet}{
  \DeltaDecl{Mengen}{A \dsep B\dsep X \dsep b\dsep c\dsep m\dsep n\dsep y}
  \DeltaDecl{Funktion}{G}
  \DeltaRow{Relationsoperatoren}{\preccurlyeq}
  \DeltaRow{Faserminimenge}{L_{G,\preccurlyeq}}
    [\FormulaRefAuto{FiberMinSet}]
}
\begin{tabproofwide}
  \proofstepwidestar[1]{\WellOrd{B,\preccurlyeq}}{\rA}
  \proofstepwidestar[2]{G\colon B\sur A}{\rA}
  \proofstepwidestar[2]{G\colon B\to A}{%
    \FormulaRefAuto{Surjektive Funktion}[def]{2}}
  \proofstepwidestar[2]{\Fib_G\colon A\to\powerset(B)}{%
    \FormulaRefAuto{\Fib_G\colon A\to \powerset(B)}{3}}
  \proofstepwidestar[5]{X\in\Fib_G[A]}{\rA}
  \proofstepwidestar[2]{\Fib_G[A]\subseteq\powerset(B)}{%
    \FormulaRefAuto{F\colon A\to B \vdash F[A]\subseteq B}{4}}
  \proofstepwidestar[2,5]{X\in\powerset(B)}{%
    \FormulaRefAuto{A\subseteq B\dsep x\in A\vdash x\in B}{6,5}}
  \proofstepwidestar[2,5]{X\subseteq B}{%
    \FormulaRefAuto{x \in \powerset(A)\eqvdash x \subseteq A}{7}}
  \proofstepwidestar[2,5]{X\neq\varnothing}{%
    \FormulaRefAuto{X\in \Fib_F[B] \vdash X\neq\varnothing}{2,5}}
  \proofstepwidestar[1,2,5]{%
    \exists m\in X\,\forall x\in X\,(m\preccurlyeq x)}{%
    \FormulaRefAuto{WellOrderPrinciple}[ax]{8,9}}

  \proofstepwidestar[11]{%
    m\in X\land\forall x\in X\,(m\preccurlyeq x)}{\rA}
  \proofstepwidestar[11]{m\in X}{\rAEa{11}}
  \proofstepwidestar[11]{%
    \forall x\in X\,(m\preccurlyeq x)}{\rAEb{11}}
  \proofstepwidestar[2,5]{%
    \exists y\in A\,X=\Fib_G(y)}{%
    \FormulaRefAuto{F\colon A\to B\dsep y\in F[A]
      \vdash \exists x\in A\,y=F(x)}{4,5}}
  \proofstepwidestar[15]{y\in A\land X=\Fib_G(y)}{\rA}
  \proofstepwidestar[15]{y\in A}{\rAEa{15}}
  \proofstepwidestar[15]{X=\Fib_G(y)}{\rAEb{15}}
  \proofstepwidestar[2,5,11]{m\in B}{%
    \FormulaRefAuto{A\subseteq B\dsep x\in A\vdash x\in B}{8,12}}
  \proofstepwidestar[11,15]{m\in\Fib_G(y)}{\rIE{17,12}}
  \proofstepwidestar[2,11,15]{G(m)=y}{%
    \FormulaRefAuto{y\in B\dsep x\in \Fib_F(y) \vdash F(x)=y}{%
      3,16,19}}

  \proofstepwidestar[21]{c\in B}{\rA}
  \proofstepwidestar[22]{G(c)=G(m)}{\rA}
  \proofstepwidestar[2,11,15,22]{G(c)=y}{%
    \FormulaRefAuto{a=b,\,b=c\vdash a=c}{22,20}}
  \proofstepwidestar[2,11,15,21,22]{c\in\Fib_G(y)}{%
    \FormulaRefAuto{y\in B\dsep x\in A\dsep F(x)=y
      \vdash x\in \Fib_F(y)}{3,16,21,23}}
  \proofstepwidestar[2,5,11,15,21,22]{c\in X}{\rIE{17,24}}
  \proofstepwidestar[2,5,11,15,21,22]{m\preccurlyeq c}{%
    \rRE{25,\rUE{13}}}
  \proofstepwidestar[2,5,11,15,21]{%
    G(c)=G(m)\rightarrow m\preccurlyeq c}{\rRI{22,26}}
  \proofstepwidestar[2,5,11,15]{%
    \forall c\in B\,
      \bigl(G(c)=G(m)\rightarrow m\preccurlyeq c\bigr)}{%
    \rUI{\rRI{21,27}}}
  \proofstepwidestar[2,5,11,15]{%
    m\in B\land
      \forall c\in B\,
        \bigl(G(c)=G(m)\rightarrow m\preccurlyeq c\bigr)}{%
    \rAI{18,28}}
  \proofstepwidestar[2,5,11,15]{%
    m\in\{\,b\in B\mid
      \forall c\in B\,
        (G(c)=G(b)\rightarrow b\preccurlyeq c)\,\}}{%
    \FormulaRefAuto{x \in \{u \in A \mid P(u)\}
      \eqvdash x \in A \land P(x)}{29}}
  \proofstepwidestar[2,5,11,15]{m\in L_{G,\preccurlyeq}}{%
    \rIE{\FormulaRefAuto{FiberMinSet}[def],30}}
  \proofstepwidestar[2,5,11,15]{%
    m\in X\land m\in L_{G,\preccurlyeq}}{\rAI{12,31}}
  \proofstepwidestar[2,5,11,15]{%
    \exists b\,(b\in X\land b\in L_{G,\preccurlyeq})}{\rEI{32}}

  \proofstepwidestar[34]{%
    n\in X\land n\in L_{G,\preccurlyeq}}{\rA}
  \proofstepwidestar[34]{n\in X}{\rAEa{34}}
  \proofstepwidestar[34]{n\in L_{G,\preccurlyeq}}{\rAEb{34}}
  \proofstepwidestar[2,5,34]{n\in B}{%
    \FormulaRefAuto{A\subseteq B\dsep x\in A\vdash x\in B}{8,35}}
  \proofstepwidestar[15,34]{n\in\Fib_G(y)}{\rIE{17,35}}
  \proofstepwidestar[2,15,34]{G(n)=y}{%
    \FormulaRefAuto{y\in B\dsep x\in \Fib_F(y) \vdash F(x)=y}{%
      3,16,38}}
  \proofstepwidestar[34]{%
    n\in\{\,b\in B\mid
      \forall c\in B\,
        (G(c)=G(b)\rightarrow b\preccurlyeq c)\,\}}{%
    \rIE{\FormulaRefAuto{FiberMinSet}[def],36}}
  \proofstepwidestar[34]{%
    n\in B\land
      \forall c\in B\,
        \bigl(G(c)=G(n)\rightarrow n\preccurlyeq c\bigr)}{%
    \FormulaRefAuto{x \in \{u \in A \mid P(u)\}
      \eqvdash x \in A \land P(x)}{40}}
  \proofstepwidestar[34]{%
    \forall c\in B\,
      \bigl(G(c)=G(n)\rightarrow n\preccurlyeq c\bigr)}{\rAEb{41}}
  \proofstepwidestar[2,11,15,34]{G(m)=G(n)}{%
    \FormulaRefAuto{a=b,\,c=b\vdash a=c}{20,39}}
  \proofstepwidestar[2,5,11,15,34]{%
    G(m)=G(n)\rightarrow n\preccurlyeq m}{%
    \rRE{18,\rUE{42}}}
  \proofstepwidestar[2,5,11,15,34]{n\preccurlyeq m}{\rRE{43,44}}
  \proofstepwidestar[5,11,34]{m\preccurlyeq n}{%
    \rRE{35,\rUE{13}}}
  \proofstepwidestar[1,2,5,11,15,34]{m=n}{%
    \FormulaRefAuto{x\leq y\dsep y\leq x\vdash x=y}[ax]{46,45}}
  \proofstepwidestar[1,2,5,11,15]{%
    \exists! b\,(b\in X\land b\in L_{G,\preccurlyeq})}{%
    \UEI{33,32,34,47}}
  \proofstepwidestar[1,2,5,11]{%
    \exists! b\,(b\in X\land b\in L_{G,\preccurlyeq})}{%
    \rEE{14,15,48}}
  \proofstepwidestar[1,2,5]{%
    \exists! b\,(b\in X\land b\in L_{G,\preccurlyeq})}{%
    \rEE{10,11,49}}
  \proofstepwidestar[1,2]{%
    \forall X\in\Fib_G[A]\,
      \exists! b\,(b\in X\land b\in L_{G,\preccurlyeq})}{%
    \rUI{\rRI{5,50}}}
\end{tabproofwide}

\FormulaThmDeltaK[Faserminima einer Wohlordnung liefern eine Rechtsinverse]{%
  \WellOrd{B,\preccurlyeq}\dsep G\colon B\sur A
  \vdash \exists F\colon A\to B\,G\circ F=\Id_A%
}{WellOrdFibersSection}{
  \DeltaDecl{Mengen}{A \dsep B}
  \DeltaDecl{Funktionen}{F \dsep G}
  \DeltaRow{Relationsoperatoren}{\preccurlyeq}
}
\begin{tabproof}
  \proofstep{1}{\WellOrd{B,\preccurlyeq}}{\rA}
  \proofstep{2}{G\colon B\sur A}{\rA}
  \proofstep{1,2}{%
    \forall X\in \Fib_G[A]\,
      \exists! b\,(b\in X\land b\in L_{G,\preccurlyeq})%
  }{\FormulaRefAuto{WellOrdFiberChoiceSet}{1,2}}
  \proofstep{1,2}{\exists F\colon A\to B\,G\circ F=\Id_A}{%
    \FormulaRefAuto{ChoiceSetRightInverse}{3}}
\end{tabproof}

\FormulaThmDeltaK[Wohlordnungssatz impliziert Auswahlaxiom]{%
  \WOP\vdash \ACsur%
}{WOPToACsur}{
  \DeltaDecl{Mengen}{A \dsep B}
  \DeltaDecl{Funktionen}{F \dsep G}
}
\begin{tabproof}
  \proofstep{1}{\WOP}{\rA}
  \proofstep{1}{\forall M\,\WellOrdable{M}}{\rIE{\FormulaRefAuto{WOP},1}}
  \proofstep{3}{G\colon B\sur A}{\rA}
  \proofstep{1}{\WellOrdable{B}}{\rUE{2}}
  \proofstep{1}{%
    \exists\,\preccurlyeq\;\WellOrd{B,\preccurlyeq}}{%
    \rIE{\FormulaRefAuto{WellOrdable},4}}

  \proofstep{6}{\WellOrd{B,\preccurlyeq}}{\rA}
  \proofstep{3,6}{\exists F\colon A\to B\,G\circ F=\Id_A}{%
    \FormulaRefAuto{WellOrdFibersSection}{6,3}}
  \proofstep{1,3}{\exists F\colon A\to B\,G\circ F=\Id_A}{\rEE{5,6,7}}

  \proofstep{1}{%
    G\colon B\sur A\rightarrow
    \exists F\colon A\to B\,G\circ F=\Id_A%
  }{\rRI{3,8}}
  \proofstep{1}{%
    \forall G\,\bigl(G\colon B\sur A\rightarrow
    \exists F\colon A\to B\,G\circ F=\Id_A\bigr)%
  }{\rUI{9}}
  \proofstep{1}{%
    \forall B\,\forall G\,\bigl(G\colon B\sur A\rightarrow
    \exists F\colon A\to B\,G\circ F=\Id_A\bigr)%
  }{\rUI{10}}
  \proofstep{1}{%
    \forall A\,\forall B\,\forall G\,\bigl(G\colon B\sur A\rightarrow
    \exists F\colon A\to B\,G\circ F=\Id_A\bigr)%
  }{\rUI{11}}
  \proofstep{1}{\ACsur}{\rIE{\FormulaRefAuto{ACsur},12}}
\end{tabproof}

\FormulaThmDeltaK[Äquivalenz von Auswahlaxiom und Wohlordnungssatz]{%
  \ACsur\eqvdash \WOP%
}{ACsurEquivWOP}{
  \DeltaDecl{Mengen}{A \dsep B}
  \DeltaDecl{Funktionen}{F \dsep G}
}
\begin{tabproofsplit}
  \proofpart{\(\vdash\)}
    \proofstep{1}{\ACsur}{\rA}
    \proofstep{1}{\WOP}{\FormulaRefAuto{ACsurToWOP}{1}}
  \closeproofpart

  \proofpart{\(\dashv\)}
    \proofstep{1}{\WOP}{\rA}
    \proofstep{1}{\ACsur}{\FormulaRefAuto{WOPToACsur}{1}}
  \closeproofpart
\end{tabproofsplit}

\chapter{Die natürlichen Zahlen als wohlgeordnetes Beispiel}

Band~10 entwickelt die Relation \(\leq_{\mathbb N}\) konkret und beweist das
Wohlordnungsprinzip der natürlichen Zahlen. Band~15 fasst ihre
Ordnungsgrundgesetze als totale Ordnung zusammen; damit folgt nun die
Klassifikation als Wohlordnung.
\section{Wohlordnung}

\FormulaThmDeltaK[Wohlordnung der natürlichen Zahlen]{%
  \WellOrd{\mathbb{N},\leq_{\mathbb{N}}}%
}{NaturalNumbersWellOrder}{%
  \DeltaRow{Relationsoperatoren}{\leq_{\mathbb{N}}}
  \DeltaRow{\textbf{Begründung}}{}
  \DeltaRow{Totale Ordnung}{\TotOrd{\mathbb{N},\leq_{\mathbb{N}}}}
    [\FormulaRefAuto{NaturalNumbersTotalOrder}[thm]]
  \DeltaRow{Wohlordnungsprinzip}{%
    B\subseteq\mathbb{N}\dsep B\neq\varnothing
    \vdash \exists a\in B\,\forall b\in B\,a\leq b}
    [\FormulaRefAuto{NaturalWellOrderingPrinciple}[thm]]
}{}
\begin{tabproof}
  \proofstep{}{\TotOrd{\mathbb{N},\leq_{\mathbb{N}}}}{%
    \FormulaRefAuto{NaturalNumbersTotalOrder}[thm]}
  \proofstep{2}{B\subseteq\mathbb{N}}{\rA}
  \proofstep{3}{B\neq\varnothing}{\rA}
  \proofstep{2,3}{\exists a\in B\,\forall b\in B\,a\leq b}{%
    \FormulaRefAuto{NaturalWellOrderingPrinciple}[thm]{2,3}}
  \proofstep{}{\WellOrd{\mathbb{N},\leq_{\mathbb{N}}}}{%
    \FormulaRefAuto{Wohlordnung}[def]{1,4}}
\end{tabproof}

\chapter{Wegweiser zu Ordnungsbeispielen}

Die Bände~12--16 entwickeln die allgemeine Ordnungstheorie von Halbordnungen
über Schranken und Paarfälle bis zu totalen Ordnungen und Wohlordnungen.
Band~10 konstruiert die Ordnung der natürlichen Zahlen; Band~17 erweitert sie
auf die ganzen Zahlen, die Bände~18 und~19 auf die rationalen und reellen
Zahlen. Band~20 behandelt endliche Mengen. Band~34 gewinnt aus
Halbverbandsoperationen Halbordnungen und verwendet dabei insbesondere die
Paarinfima und Paarsuprema aus Band~14.

\end{document}
